\documentclass[openany]{book}

\usepackage[margin=1in]{geometry}
\usepackage{amsmath,amsfonts,amsthm, amssymb}
\usepackage{yhmath}
\usepackage{mathrsfs}
\usepackage{mathtools}
\usepackage{xcolor}
\usepackage{graphicx}
\usepackage{comment}
\usepackage{tikz-cd}
\usepackage{quiver}
\usepackage{hyperref,cleveref}
\renewcommand{\familydefault}{ppl}
\newcommand{\curl}{\text{curl}}
\newcommand{\diverg}{\text{div}}
\newcommand{\tr}{\text{tr}}
\newcommand{\R}{\mathbb{R}}
\newcommand{\E}{\mathbb{E}}
\newcommand{\Z}{\mathbb{Z}}
\newcommand{\C}{\mathbb{C}}
\newcommand{\F}{\mathbb{F}}
\newcommand{\la}{\langle}
\newcommand{\ra}{\rangle}
\newcommand{\colim}{\text{colim}}
\DeclareMathOperator{\im}{im}
\DeclareMathOperator{\disc}{disc}
\let\oldemptyset\emptyset
\let\emptyset\varnothing
\newcommand{\tor}{\text{Tor}}
\newcommand{\id}{\text{id}}
\newcommand{\ext}{\text{Ext}}
\newcommand{\ptop}{\text{PTop}}
\newcommand{\pt}{\text{pt}}
\newcommand{\ach}{\text{Ach}}
\newcommand{\Q}{\mathbb{Q}}
\newcommand{\gal}{\text{Gal}}
\newcommand{\fraccomma}{\genfrac{}{}{0pt}{}{}{,}}
\definecolor{wikipediadarkblue}{rgb}{0.023, 0.270, 0.676}
\hypersetup{
    colorlinks,
    citecolor=black,
    filecolor=black,
    linkcolor=blue,
    urlcolor=red
}

% \hypersetup{
%     urlbordercolor={1 0 0}, % Red box for URLs
%     linkbordercolor=blue, % Blue box for internal links
%     % the values are in RGB format from 0 to 1
% }


\input{hui_r.tex}

\title{Complex Analysis Qualifying Exam Review
\\ 
\vspace{0.4cm}
\large Spring 2026}



% \author{(Please email hsun95@jh.edu if you see typos)}
% \date{\today}


\begin{document}

\maketitle

\tableofcontents
\newpage



\chapter{Basic calculus in the complex domain}

\begin{prop}[Cauchy-Riemann Equations]
    If $f:\Omega\to\C$ is holomorphic, then $\frac{\partial f}{\partial x}, \frac{\partial f}{\partial y}$ exist and are continuous on $\Omega$, and 
    \begin{equation*}
        \frac{\partial f}{\partial x}=\frac{1}{i}\frac{\partial f}{\partial y}
    \end{equation*}
    If we write $f=u+iv$, then the above is equivalent to 
    \begin{equation*}
        \begin{cases}
            \frac{\partial u}{\partial x}=\frac{\partial v}{\partial y}\\
            \frac{\partial v}{\partial x}=-\frac{\partial u}{\partial y}
        \end{cases}
    \end{equation*}
\end{prop}

\begin{prop}[Converse of Cauchy-Riemann]
    If $f:\Omega\to\C$ is $C^1$ and satisfies Cauchy-Riemann equations, then $f$ is holomorphic.
\end{prop}

\begin{prop}
    If $f:\overline{\Omega}\to\C$ is holomorphic, where $z_0\in\Omega$ and $D_r(z_0)\subset\Omega$, then for $z\in D_r(z_0)$, $f$ is given by a convergent power series 
    \begin{equation*}
        f(z)=\sum_{k=0}^\infty a_k(z-z_0)^k
    \end{equation*}
    where 
    \begin{equation*}
        a_k=\frac{1}{2\pi i}\int_{\partial\Omega}\frac{f(\zeta)}{(\zeta-z_0)^{k+1}}d\zeta=\frac{f^{(k)}(z_0)}{k!}
    \end{equation*}
\end{prop}
(This follows from Cauchy integral formula).


% \begin{defn}[complex log]
%     The logorithmic function extended to the complex plane is $\log:\C\setminus\R^-\to\C$, given by 
%     \begin{equation*}
%         \log z=\frac{1}{z}\frac{1}{\zeta}d\zeta
%     \end{equation*}
% \end{defn}


\begin{thm}
    Let $f:\Omega\to\C$ be holomorphic, where $\Omega$ is open. Let $p\in\Omega$, assume $f'(p)\neq 0$. Then there exists a neighborhood $O$ of $p$ and a neighborhood $U$ of $q=f(p)$ such that $f:O\to U$ is bijective, where the inverse $g=f^{-1}: U\to O$ is holomorphic. For $z\in O, w=f(z)$, we have 
    \begin{equation*}
        g'(w)=\frac{1}{f'(z)}
    \end{equation*}
\end{thm}


\begin{prop}
    Suppose $\Omega\subset\C$ is convex. Assume $f$ is holomorphic in $\Omega$ and there exists $a\in\C$ such that 
    \begin{equation*}
        \text{Re } af'>0 \text{ on } \Omega
    \end{equation*}
    then $f$ is one-to-one onto its image $f(\Omega)$.
\end{prop}


\begin{defn}
    square roots, complex logs
\end{defn}


\chapter{Cauchy integral theorem}









\chapter{Useful facts}
\begin{enumerate}
    \item $e^z$ has the following power series expansion: 
    \begin{equation*}
        e^z=\sum_{j=0}^\infty\frac{z^j}{j!}
    \end{equation*}
    \item $\cos, \sin$ has the following power series expansions:
    \begin{equation*}
        \cos z=
    \end{equation*}
    \begin{equation*}
        \sin z=
    \end{equation*}
\end{enumerate}






\end{document}